\clearpage
\fancyhead[L]{\small TOBIAS MAYER. ASTRONOMICAL OBSERVATIONS.}
\hypertarget{native-649-lower}{}
\begin{center}\small Göttingische gelehrte Anzeigen. No. 160. pp. 1585..1588. 6 October 1827.\\[3pt]\rule{.72\linewidth}{.3pt}\end{center}

\emph{Astronomical observations made at Göttingen from \textup{1756} to \textup{1761} by \textup{\textsc{Tobias Mayer}}. In two parts. Published by order of the commissioners of latitude.} London 1826. The first part has 92 pages and the second 62 pages, both in folio.

After seventy years, the observations of the great Göttingen astronomer \textsc{Tobias Mayer} have at last been accorded the distinction of being printed at the expense of a foreign authority that gladly uses the considerable resources at its disposal for the benefit of astronomy. The manuscript of the observations from the years 1757..1761 had long since come into the possession of Mr.~von \textsc{Zach}, who had shared it with the English commission for longitude at sea. The manuscript of the observations from the first year was still in the archives of the local Royal Society and was lent for the intended printing.

By far the greater part of the observations was made with \textsc{Bird}'s mural quadrant. \textsc{Mayer} had directed his chief attention to compiling the catalogue of zodiacal stars, and observations of these stars therefore constitute the main content of the journal. They begin on 8 February 1756. At first, up to 23 May, there are only the transit times, already reduced to the middle wire. From 23 May onward, by contrast, the observations appear in full in their original form. The observations of 1756 fill the whole first part. Those for 1757 occupy 35 pages in the second part, followed by 6 pages up to the middle of 1758. At this point \textsc{Mayer}'s observations with the mural quadrant practically cease altogether. On 5 April 1759 he presented his catalogue of zodiacal stars to the local Royal Society. From the first days of May 1759, observations of \textsc{Halley}'s comet with an equatorially mounted telescope. From June and July 1760, observations with the mural quadrant by \textsc{Mayer}'s worthy pupil \textsc{Niebuhr}. Finally, from 6 June 1761, the observation of the transit of Venus across the Sun. Besides the observations of zodiacal stars, in the early years we find very many observations of the Sun, several of Venus and the Moon, and a few of Mars, Jupiter and Saturn, as well as various stellar occultations.

From these contents one can see that a fresh analysis of these observations from the ground up could not yield results as extensive as those obtained from the contemporary observations of \textsc{Bradley} through \textsc{Bessel}'s excellent analysis. \textsc{Mayer} had set himself a limited \SourceLeaf{650}main objective, the compilation of a catalogue of zodiacal stars, and he himself based that catalogue on his observations. It must also be borne in mind that the scantiness of \textsc{Mayer}'s resources poses major obstacles to reducing the observations as absolute observations and that we now lack knowledge of some of the ancillary aids he also used. \textsc{Mayer} had no transit telescope. Reducing the right ascensions observed with the mural quadrant requires knowledge of the instrument's irregular departures from the meridian. For this purpose \textsc{Mayer} also used corresponding altitudes of the Sun taken with a small movable quadrant. We find such readings preserved, however, from only two days. Similar difficulties arise with the zenith distances. \textsc{Mayer} had no zenith sector. Reliable knowledge of the collimation error could therefore be obtained only by rehanging the instrument, and this was done only once. All the zenith distances observed before the rehanging are made still less reliable by \textsc{Mayer}'s own remark that the central pivot was found to be loose. Several other remarks could be added, but there is no place for them here.

On the other hand, individual cases may still often arise in which astronomers will be very glad to be able to use the original observations, even if only as differential observations. As is well known, the manuscript has already on one occasion rendered the great service of supplying an early observation of Uranus, which \textsc{Mayer} unknowingly observed as a fixed star on 25 Sept.~1756.

The introduction also reprints two papers from the Opera inedita, namely \textsc{Mayer}'s lecture Observationes astronomicae quadrante murali habitae and \textsc{Lichtenberg}'s remarks on the thermometer used by \textsc{Mayer}. At the end there is a list, submitted by Prof.~\textsc{Bessel}, of misprints in \textsc{Bradley}'s observations.
